% !TEX root = ../main.tex
\section{Инструменты и методы реализации}
\addcontentsline{toc}{section}{Инструменты и методы реализации}

\subsection{Технологический стек}
Для реализации программного обеспечения был выбран язык программирования Python ввиду его богатой экосистемы библиотек для компьютерного зрения, обработки данных и создания графических интерфейсов. Основные компоненты технологического стека включают:

\begin{itemize}
    \item \textbf{YOLOv8} --- современная архитектура свёрточных нейронных сетей для задач детекции объектов в реальном времени. Была использована предобученная модель, доступная через платформу Roboflow, обученная на датасете объёмом более 21\,000 изображений игральных карт. Данная модель послужила базой для дальнейшей донастройки (fine-tuning) под специфику интерфейса целевого покер-рума.
    
    \item \textbf{PyQt6} --- фреймворк для разработки кроссплатформенных графических интерфейсов. Используется для создания Heads-Up Display (HUD) --- оверлейного окна, отображающего рекомендации поверх клиентского приложения покер-рума без нарушения его работы.
    
    \item \textbf{OpenCV} --- библиотека компьютерного зрения, применяемая для предварительной обработки изображений (масштабирование, нормализация яркости, фильтрация шума) перед подачей в нейросеть.
    
    \item \textbf{NumPy, Pandas} --- инструменты для эффективной обработки и хранения данных о диапазонах рук, статистики игровых ситуаций и результатов распознавания.
\end{itemize}

\subsection{Подготовка датасета и обучение модели}
Одним из ключевых этапов разработки стал процесс создания специализированного датасета для распознавания элементов интерфейса покер-рума \textit{ReplayPoker}. Учитывая отсутствие готовых размеченных данных для данной платформы, был реализован следующий подход:

\begin{enumerate}
    \item Сбор исходных изображений игровых столов путём ручного скриншотирования различных игровых ситуаций (различные позиции, стеки, этапы торговли).
    
    \item Ручная разметка изображений с использованием инструментов аннотации (bounding boxes для карт, кнопок действий, индикаторов стеков и позиций).
    
    \item Аугментация данных для повышения устойчивости модели к вариациям освещения, масштаба и незначительным изменениям интерфейса.
    
    \item Дообучение предобученной модели YOLOv8 на подготовленном датасете с подбором оптимальных гиперпараметров (learning rate, batch size, количество эпох) для достижения баланса между скоростью инференса и точностью детекции.
\end{enumerate}

Данный подход позволил адаптировать общую модель распознавания карт под специфику конкретного покер-рума, обеспечив высокую точность определения игровых элементов при минимальной задержке обработки.

\subsection{Архитектура парсера игрового стола}
Разрабатываемый парсер игровой ситуации реализован как модульное приложение, выполняющее следующие функции:

\begin{itemize}
    \item Захват изображения области экрана, занятой клиентом покер-рума.
    
    \item Предварительная обработка изображения и последующая детекция ключевых элементов интерфейса с помощью дообученной модели YOLOv8.
    
    \item Извлечение семантической информации: определение карт игрока и борда, позиции за столом, размеров стеков и текущих ставок.
    
    \item Формирование структурированного представления игровой ситуации для передачи в модуль принятия решений.
    
    \item Отображение результатов анализа через оверлейный HUD на базе PyQt6, не взаимодействующий напрямую с окном покер-клиента (во избежание нарушения правил платформы).
\end{itemize}
